%% RescueSched --- INFOCOM submission
%% Build: latexmk -pdf main.tex   (uses pdflatex + bibtex + IEEEtran.bst)
%% Class: IEEEtran conference (double-column, 10pt, letterpaper)
\documentclass[conference]{IEEEtran}

% --- IEEE requires these to be loaded ---
\usepackage{cite}
\usepackage{amsmath,amssymb,amsfonts}
\usepackage{algorithmic}
\usepackage{algorithm}
\usepackage{graphicx}
\usepackage{textcomp}
\usepackage{xcolor}
\usepackage{booktabs}
\usepackage{array}
\usepackage{url}

% Correct bad hyphenation here
\hyphenation{op-tical net-works semi-conduc-tor}

\graphicspath{{figures/}}

% Editorial note macro (turn off before submission by setting to {})
\newcommand{\note}[1]{\textcolor{red}{[#1]}}

\begin{document}

\title{RescueSched: Task-Level SLO Rescue via Selective Descriptor
Migration for Microsecond-Scale RPC Services}

% INFOCOM is single-blind: author block is shown at submission.
% Comment out and use \author{Anonymous} if a given year is double-blind.
\author{%
\IEEEauthorblockN{First Author\IEEEauthorrefmark{1},
Second Author\IEEEauthorrefmark{2}}
\IEEEauthorblockA{\IEEEauthorrefmark{1}Affiliation One,
\texttt{email1@example.com}}
\IEEEauthorblockA{\IEEEauthorrefmark{2}Affiliation Two,
\texttt{email2@example.com}}
}

\maketitle

\begin{abstract}
Latency-critical RPC services on commodity multi-core servers commonly
pin one worker per core with a per-core request queue. Under bursty,
skewed, or heavy-tailed load, transient sper-core queue imbalance causes
P99/P99.9 tail-latency SLO violations even when the server is not
globally overloaded. Existing work stealing, work shedding, and
threshold-based migration answer only \emph{which core has too much
work}, not \emph{which queued request will miss its SLO}, \emph{whether
it can still be saved after migration}, and \emph{whether a target core
can absorb it safely}. We present RescueSched, which reframes intra-host
migration from load balancing to \emph{SLO rescue}: it migrates only a
queued-but-not-running request that is locally doomed, remotely
rescuable under a migration-cost-aware completion-time estimate, and
safe for the target (no predicted increase in target-side violations).
On a discrete-event simulator of a 16-core host, RescueSched reduces the
SLO violation ratio on a heavy-tailed workload at $\rho{=}0.85$ from
$0.20$ (work stealing) and $0.15$ (proactive migration) to $0.094$,
while migrating \emph{fewer} requests. Ablations show the gain comes
from a higher beneficial-migration ratio rather than migration volume,
and that RescueSched correctly suppresses migration under global
overload.
\end{abstract}

\begin{IEEEkeywords}
tail latency, SLO, microsecond-scale RPC, request migration, multi-core
scheduling
\end{IEEEkeywords}

\section{Introduction}
\label{sec:intro}
% !TEX root = ../main.tex
Modern latency-critical RPC services often run on commodity multi-core
servers with one worker pinned to each CPU core. To avoid centralized
queue contention and preserve cache locality, each worker maintains a
local request queue, and incoming packets are distributed by RSS, flow
hashing, or a user-space ingress thread. This design scales well under
balanced load, but it can create transient per-core queue imbalance
under bursty arrivals, skewed flow distribution, and heterogeneous
service times. For microsecond-scale RPC, even a short period of local
queue buildup can cause P99 or P99.9 tail-latency SLO violations.

Existing work stealing improves work conservation by letting idle cores
pull tasks from busy queues; work shedding and threshold-based migration
react when a queue becomes overloaded. However, these mechanisms mainly
answer one question: \emph{which core has too much work?} They do not
explicitly answer three finer-grained SLO questions:
(i) which queued request will miss its SLO if it stays local;
(ii) which of those requests can still meet its SLO after migration;
(iii) which target core can accept the migrated request without becoming
risky itself.

\note{Sharpen the gap statement and add one concrete motivating number
from the evaluation here.}

\textbf{Contributions.} This paper makes three contributions:
\begin{itemize}
  \item We define \emph{task-level rescuability} for microsecond-scale
  RPC scheduling, using local-vs-remote completion-time counterfactuals
  under a migration-cost model (\S\ref{sec:model}).
  \item We design \emph{migration-safe target selection}, which prevents
  cross-core migration from transferring SLO risk from the source core
  to the target core (\S\ref{sec:design}).
  \item We implement RescueSched as a commodity user-space prototype and
  show that selective descriptor migration reduces SLO violations and
  useless migrations compared with work stealing, work shedding, and
  threshold-based migration baselines (\S\ref{sec:eval}).
\end{itemize}

\section{Background and Motivation}
\label{sec:background}
% !TEX root = ../main.tex
\note{Background: microsecond-scale RPC tail latency, per-core queues,
d-FCFS vs single-queue, and why queue length alone is a poor migration
signal.} Prior microsecond-scale schedulers~\cite{placeholder-shinjuku}
and work-conserving stealing~\cite{placeholder-workstealing} target work
conservation rather than per-request SLO outcomes.

RescueSched is not a generic load balancer. It does not aim to equalize
queue lengths. Instead, it treats cross-core migration as a scarce and
costly SLO-rescue action: a migration is useful only if it changes the
SLO outcome of a request without creating new target-side violations.
The central problem is:

\begin{quote}
Given a commodity multi-core RPC server with per-core request queues and
no preemption of running requests, can we reduce the SLO violation ratio
by selectively migrating only queued-but-not-running requests that are
locally doomed, remotely rescuable, and safe for the target core?
\end{quote}

\textbf{Scope and non-goals.} We target a single latency-critical RPC
service on one multi-core server with $N$ pinned worker cores, one FIFO
queue per core, microsecond-scale service times, request-level SLOs, no
preemption, no kernel modification, and no SmartNIC/programmable-switch
dependence. The scheduler migrates only queued request descriptors; it
never migrates running requests, application state, or payloads. We do
not attempt to solve global overload: if all cores are saturated with no
slack, RescueSched suppresses migration.

\section{System Model}
\label{sec:model}
% !TEX root = ../main.tex
\subsection{Cores, Queues, and Requests}
Let the server have cores $\mathcal{C}=\{1,\dots,N\}$. Each core $c$
maintains a FIFO queue $Q_c(t)$ and may execute one request with
estimated remaining time $R_c(t)$ ($R_c(t){=}0$ if idle). The estimated
queued workload is
\begin{equation}
W_c(t) = R_c(t) + \sum_{r_i \in Q_c(t)} \hat{s}_i ,
\end{equation}
where $\hat{s}_i$ is the estimated service time of $r_i$. Each request is
$r_i=\langle id_i,a_i,D_i,m_i,\hat{s}_i,c_i,state_i\rangle$ with deadline
$D_i=a_i+SLO_i$. The scheduler considers only requests with
$state_i=\texttt{queued}$; running requests are never migrated.

\subsection{Service-Time Estimation}
Per-method statistics maintain an EWMA estimate
$\mu_m(t)=(1-\alpha)\mu_m(t-1)+\alpha s_i$, so $\hat{s}_i=\mu_{m_i}(t)$.
For heavy-tailed methods a conservative quantile $\hat{s}_i=Q_p(S_{m_i})$
(e.g.\ $p{=}0.90$) may be used as an ablation.

\subsection{Local and Remote Completion-Time Counterfactuals}
For $r_i$ at source core $s$ with $Before(i,s)$ the requests ahead of it,
\begin{equation}
\hat{C}_{\text{local}}(i,s)= t + R_s(t)
  + \!\!\sum_{r_k \in Before(i,s)}\!\! \hat{s}_k + \hat{s}_i .
\end{equation}
Request $r_i$ is \emph{locally doomed} if
$\hat{C}_{\text{local}}(i,s) > D_i$. With migration cost
$M_{s,j}(i)=M_{desc}+M_{sync}+M_{numa}(s,j)$ (same-NUMA
$\Rightarrow M_{numa}{=}0$), appending $r_i$ to the tail of $Q_j$ gives
\begin{equation}
\hat{C}_{\text{remote}}(i,j)= t + M_{s,j}(i) + R_j(t)
  + \!\!\sum_{r_k \in Q_j(t)}\!\! \hat{s}_k + \hat{s}_i .
\end{equation}
$r_i$ is \emph{remotely feasible} on $j$ if
$\hat{C}_{\text{remote}}(i,j)+\epsilon \le D_i$, where $\epsilon$ absorbs
estimation and cost uncertainty.

\subsection{Rescuability and Target Safety}
$r_i$ is \emph{rescuable} from $s$ to $j$ iff
\begin{equation}
\hat{C}_{\text{local}}(i,s) > D_i
\;\wedge\;
\hat{C}_{\text{remote}}(i,j)+\epsilon \le D_i .
\end{equation}
This excludes \emph{needless} migration (would meet SLO locally) and
\emph{unsaved} migration (would miss SLO anyway). The remote slack is
$Slack_{\text{remote}}(i,j)=D_i-\hat{C}_{\text{remote}}(i,j)$.
A target is \emph{migration-safe} if the incremental predicted risk
$\Delta Risk(j|i,t)=Risk_{\text{after}}(j|i,t)-Risk(j,t)\le\theta$,
where $Risk(j,t)=\sum_{r_x\in Q_j(t)}\mathbb{I}[\hat{C}_{\text{local}}
(x,j)>D_x]$ and the default strict setting is $\theta=0$.

\section{RescueSched Design}
\label{sec:design}
% !TEX root = ../main.tex
\subsection{Design Insight}
Queue length identifies overloaded cores but not which task is worth
rescuing or which target can absorb it safely. RescueSched therefore uses
queue length only as a risk trigger; the final decision depends on the
local counterfactual, the remote counterfactual, migration cost, remote
deadline slack, and incremental target-side risk.

\subsection{Scheduling Loop}
On each scheduling epoch (or queue-update event), RescueSched updates
per-core statistics, identifies risky source cores and low-risk targets,
and for each source scans the first $K$ \emph{locally doomed} candidates.
For each candidate it tests up to $H$ targets, keeping pairs that are
remotely feasible and target-safe, scores them, and commits at most $B$
descriptor-only migrations per epoch, updating simulated target state
between selections.

\begin{algorithm}[t]
\caption{RescueSched main loop (per epoch $t$)}
\label{alg:main}
\begin{algorithmic}[1]
\STATE update $R_c,W_c,Risk(c,t)$ for all $c$
\STATE $S_{risky}\!\gets\!\{s: Risk(s,t)\!>\!0 \lor W_s \text{ over trigger}\}$
\STATE $T_{safe}\!\gets\!\{j: Risk(j,t)\!=\!0 \lor W_j \text{ under trigger}\}$
\STATE $Pairs \gets \emptyset$
\FOR{$s \in S_{risky}$}
  \FOR{$i \in$ top-$K$ locally-doomed of $Q_s$}
    \IF{$\hat{C}_{\text{local}}(i,s)\le D_i$} \STATE \textbf{continue} \ENDIF
    \FOR{$j \in$ top-$H$ of $T_{safe}$, $j\neq s$}
      \IF{$\hat{C}_{\text{remote}}(i,j)+\epsilon > D_i$} \STATE \textbf{continue} \ENDIF
      \IF{$\Delta Risk(j|i,t) > \theta$} \STATE \textbf{continue} \ENDIF
      \STATE $Pairs \mathrel{+}= (i,j,\textsc{Score}(i,j))$
    \ENDFOR
  \ENDFOR
\ENDFOR
\STATE sort $Pairs$ by descending score; $Selected \gets \emptyset$
\FOR{$(i,j)\in Pairs$}
  \IF{$|Selected|\!\ge\!B$} \STATE \textbf{break} \ENDIF
  \IF{$i$ selected $\lor$ $j$ no longer safe} \STATE \textbf{continue} \ENDIF
  \STATE $Selected \mathrel{+}= (i,j)$; update simulated state of $j$
\ENDFOR
\STATE execute descriptor-only migrations for $Selected$
\end{algorithmic}
\end{algorithm}

\subsection{Scoring}
With $\gamma{=}0,\theta{=}0$ in the minimal version, the score reduces to
$Score(i,j)=Slack_{\text{remote}}(i,j)-\beta M_{s,j}(i)$, ranking
migrations by robustness to estimation error net of cost. The full form
adds $\gamma\cdot local\_lateness-\lambda\cdot\Delta Risk$.

\subsection{Complexity}
With per-core prefix sums, $\hat{C}_{\text{local}}$ is $O(1)$ per request
after snapshotting, and target risk is approximated in $O(KH)$ per epoch,
giving naive per-epoch cost $O(|S_{risky}|\cdot K\cdot H\cdot L)$.

\subsection{Migration Usefulness Taxonomy}
Each committed migration is labeled post-hoc as \emph{beneficial},
\emph{needless}, \emph{unsaved}, or \emph{harmful}, yielding the
beneficial-migration ratio (BMR), useless-migration ratio (UMR), and
rescue-per-migration (RPM). These measure whether the scheduler migrates
\emph{precisely}, not merely frequently.

\section{Evaluation}
\label{sec:eval}
% !TEX root = ../main.tex
\subsection{Setup}
We evaluate on a discrete-event simulator of a single host with $N{=}16$
pinned worker cores. New requests are randomly assigned to cores; only
waiting requests migrate; running requests do not. Unless noted, the
defaults are $check{=}1\,\mu s$, $K{=}16$, $H{=}4$, $B{=}1$,
$\epsilon{=}2\,\mu s$, $\theta{=}0$, and $M{=}0.5\,\mu s$. Workloads are
W1 (near-saturation), W2 (MMPP burst-skew), and W3 (heavy-tailed).
Baselines are RSS-local FIFO (d-FCFS), work stealing, intra-host
proactive migration, and threshold/shortest-queue migration. All numbers
are medians over 5 seeds unless a 10-seed CI is reported.

\note{Regenerate Table~\ref{tab:main} and figures from
artifacts/step-15-rescuesched and step-17-rescuesched-closure via
scripts/rescue\_analysis.py before submission.}

\begin{table}[t]
\centering
\caption{W3 heavy-tail, median over 5 seeds. RescueSched lowers SLO
violation while migrating fewer requests than proactive migration.}
\label{tab:main}
\small
\begin{tabular}{@{}llrrrr@{}}
\toprule
$\rho$ & Method & SLO viol. & Move rate & BMR & UMR \\
\midrule
0.85 & Work stealing        & 0.203 & 0.328 & 0.00 & 0.00 \\
0.85 & Proactive mig.\ (M0) & 0.154 & 0.400 & 0.00 & 0.00 \\
0.85 & \textbf{RescueSched} & \textbf{0.094} & 0.286 & 1.00 & 0.00 \\
\bottomrule
\end{tabular}
\end{table}

\subsection{Main Result}
On W3 at $\rho{=}0.85$ RescueSched cuts the SLO violation ratio to
$0.094$ from $0.203$ (work stealing) and $0.154$ (proactive migration),
while its move rate ($0.286$) is \emph{lower} than proactive migration's
($0.400$)---the gain is quality, not volume. Its P99/P99.9 are not always
the lowest; RescueSched optimizes SLO rescue and migration usefulness,
not raw tail latency, and we state this trade-off explicitly.

\subsection{Ablations}
Removing rescuability (\emph{NoRescuable}) migrates far more
($0.60$ vs $0.29$) but collapses BMR to $0.31$, raises UMR to $0.69$, and
degrades SLO violation to $0.169$---confirming that task-level
rescuability is necessary. Removing target safety leaves \emph{actual}
harmful migrations at $0$ under the current FIFO append-to-tail model; it
reduces only \emph{predicted} target-risk exposure. We report this as a
model boundary rather than a strong claim.

\subsection{Sensitivity and Boundaries}
$check{=}1\,\mu s$ is the operating point; $5\,\mu s$ misses the rescue
window. Budget $B{=}1$ dominates $B{=}2,4$, which add unsaved migrations.
Under W1 global overload ($\rho{=}0.95$) RescueSched performs $0$
migrations, confirming it does not fabricate capacity.

\section{Related Work}
\label{sec:related}
% !TEX root = ../main.tex
\note{Related work to cite and contrast:}
\begin{itemize}
  \item \textbf{Microsecond-scale scheduling / kernel bypass:} Shinjuku,
  Shenango, Caladan, ZygOS, IX, Arachne, Persephone~\cite{placeholder-shinjuku}.
  \item \textbf{Work stealing / conservation:} classic work stealing,
  BWS, and RPC-level stealing~\cite{placeholder-workstealing}.
  \item \textbf{Threshold / shedding migration:} ALTO-style threshold
  migration and shortest-queue placement.
  \item \textbf{Load balancing / power-of-$d$:} JSQ, power-of-two-choices,
  and their RPC variants.
\end{itemize}
RescueSched differs by treating migration as a deadline-driven rescue
action gated on local/remote completion counterfactuals and target-side
safety, rather than as queue-length equalization.

\section{Conclusion}
\label{sec:conclusion}
% !TEX root = ../main.tex
We presented RescueSched, which reframes intra-host request migration for
microsecond-scale RPC from load balancing to task-level SLO rescue. By
migrating only queued requests that are locally doomed, remotely
rescuable under a cost-aware completion estimate, and safe for the target,
RescueSched reduces SLO violations on heavy-tailed and bursty workloads
while migrating fewer requests, and it suppresses migration under global
overload. Future work includes a CloudLab user-space prototype, an
observable service-time predictor to replace oracle estimates, and a
stronger target-side harm model.

\bibliographystyle{IEEEtran}
\bibliography{refs}

\end{document}
